\section{Project Plan}

As part of my thesis research, I am investigating preference elicitation of Large Language Models (LLMs).
This is done under the lens of economics, measuring the effect of such preferences on markets.
This project applies GRUMs to LLM preference elicitation.
Unlike the Bradley-Terry (BT) model, currently used in my thesis research, which treats prompt variations as random noise, we assume these variations create specific "personas" with distinct preferences.
We will use GRUMs to model how the LLM's internal state (the "agent") interacts with the items (the "alternatives").

\subsection{Methodology}
We propose a framework where LLM prompts serve as agents (personas) and predefined token outputs (e.g., colors) as alternatives.

\begin{itemize}
    \item \textbf{Agents ($x_i$):}
    We define the "agent" as the model's internal state after processing the initial prompt. We will use the \textbf{last hidden layer} to represent this state.
    Since these vectors are high-dimensional, we will project them down to a lower dimension (using PCA or a learnable embedding layer) to make learning the interaction matrix ($B$) computationally tractable (Therefore, dimension size $K$ is a hyperparameter).
    \item \textbf{Alternatives ($z_j$):}
    The items (e.g., colors) are represented by their standard token embeddings (which might also require a dimensionality reduction) (Therefore, dimension size $L$ is also a hyperparameter).
    \item \textbf{Utility Function:}
    $u_{ij} = \delta_j x_i B z_j^{\top} + \epsilon_{ij}$
    \begin{itemize}
        \item $\delta_j$ is the "intrinsic" preference (context-free).
        \item $B$ is the interaction matrix that captures how the hidden state modifies preferences.
    \end{itemize}
\end{itemize}

\subsection{Experiment 1: Intrinsic Preference Validation}
We will train both a baseline BT model and our GRUM on the same comparison data from the Qwen-2.5 model family, and compare the learned parameters.

\textbf{Conjecture:} The GRUM's intrinsic parameters ($\delta$) will highly correlate with the BT parameters ($\beta$).
We expect BT to act as the "average" of the detailed preferences captured by GRUM.
See Fig. \ref{fig:experiments} Left.

\subsection{Experiment 2: The persona Effect}

We will analyze the interaction matrix $B$ by predicting preferences under two distinct hidden states (e.g., one generated by a "Frugal" prompt vs. a "Luxury" prompt).

\textbf{Conjecture:} The matrix $B$ will be non-zero, causing "Rank Reversals".
For example, the Frugal state might rank "Item A" high, while the Luxury state ranks it low, proving preferences depend on the hidden state.
See Fig. \ref{fig:experiments} Center.

To ensure the identified personas represent stable internal states, rather than sensitivity to specific phrasing, we will measure the similarity within a persona cluster and dissimilarity between persona clusters (i.e., each prompt induces an agent's persona $x_i$, we will look at similarity between the embeddings of prompts with similar semantic meanings such as "You prefer budget options" and "Act as a cost-conscious buyer" and also measure dissimilarity between "Frugal persona" prompts and "Luxury persona" prompts

\subsection{Experiment 3: Predictive Performance}

We will compare the Negative Log-Likelihood (NLL) of both models on a held-out test set of new prompts.

\textbf{Conjecture:} GRUM will achieve a lower NLL (better fit) than BT.
Since GRUM is condition on the hidden state, it can predict the specific variation in preference, whereas BT has to treat that variation as random error.
See Fig. \ref{fig:experiments} Right.

\begin{figure}
    \centering
    \includegraphics[width=1\linewidth]{experiments.png}
    \caption{
    \textbf{(Left)} Correlation of intrinsic preferences between GRUM and BT Baseline.
    \textbf{(Center)} Rank reversals illustrating the dependency of preferences on the model's hidden state.
    \textbf{(Right)} Superior predictive performance (lower NLL) of GRUM compared to Bradley-Terry.}
    \label{fig:experiments}
\end{figure}